Set
\[
R=[1/4,3/4]\times[0,1]\times[0,1]
\]
and let \(P,R_0,R_1\) be independent, with \(P\) uniform on \([1/4,3/4]\) and each \(R_a\) uniform on \([0,1]\).  Define the continuous map \(\phi:R\to K_{1/4}\) by
\[
\phi(P,R_0,R_1)
=\bigl((1-P)(1-R_0),(1-P)R_0,P(1-R_1),PR_1\bigr),
\]
where the coordinates are ordered as \((q_{00},q_{01},q_{10},q_{11})\), and denote the pushforward of the product uniform law under \(\phi\) by \(\nu\).  The coordinates of \(\phi\) are nonnegative, sum to one, and satisfy
\[
p(\phi(P,R_0,R_1))=P\in[1/4,3/4].
\]
Thus \(\phi(R)\subseteq K_{1/4}\) by the definition of \(K_{\epsilon}\).  Conversely, if \(q\in K_{1/4}\), then
\[
P=p(q),\qquad R_0=\frac{q_{01}}{1-p(q)},\qquad R_1=\frac{q_{11}}{p(q)}
\]
are well defined because \(1/4\le p(q)\le3/4\), belong to the respective intervals, and satisfy \(\phi(P,R_0,R_1)=q\).  Hence \(\phi(R)=K_{1/4}\).

The measure \(\nu\) has full support on \(K_{1/4}\).  Indeed, if \(G\) is a nonempty relatively open subset of \(K_{1/4}\), continuity and surjectivity of \(\phi\) imply that \(\phi^{-1}(G)\) is a nonempty relatively open subset of the full-dimensional rectangular box \(R\).  Such a set has positive three-dimensional Lebesgue measure, so \(\nu(G)>0\).

The required moments follow by direct calculation.  Independence and the uniform laws give
\[
\mathbb E[P]=\frac12,qquad
\mathbb E[P^2]=\mathbb E[(1-P)^2]=\frac{13}{48},qquad
\mathbb E[P(1-P)]=\frac{11}{48},
\]
and, for \(a\in\{0,1\}\),
\[
\mathbb E[R_a]=\mathbb E[1-R_a]=\frac12,qquad
\mathbb E[R_a^2]=\mathbb E[(1-R_a)^2]=\frac13,qquad
\mathbb E[R_a(1-R_a)]=\frac16.
\]
For a cell \((a,y)\), let \(\mathbf e_{ay}\) denote its unit count vector.  By the definition of the degree-two Bernstein polynomials, there are three exhaustive cases.  First, for every \((a,y)\),
\[
\int B_{2\mathbf e_{ay}}^{(2)}(q)\,d\nu(q)
=\mathbb E[q_{ay}^{,2}]
=\frac{13}{48}\frac13
=\frac{13}{144}.
\]
Second, for the two distinct outcomes within one treatment arm,
\[
\int B_{\mathbf e_{a0}+\mathbf e_{a1}}^{(2)}(q)\,d\nu(q)
=2\mathbb E[q_{a0}q_{a1}]
=2\frac{13}{48}\frac16
=\frac{13}{144}.
\]
Third, for \(y,z\in\{0,1\}\), independence of \(R_0\) and \(R_1\) yields
\[
\int B_{\mathbf e_{0y}+\mathbf e_{1z}}^{(2)}(q)\,d\nu(q)
=2\mathbb E[q_{0y}q_{1z}]
=2\frac{11}{48}\frac12\frac12
=\frac{11}{96}.
\]
The first two cases are exactly the count vectors \(c\in\mathcal I_2\) for which \(c_{00}+c_{01}=2\) or \(c_{10}+c_{11}=2\); the third case comprises all remaining count vectors.  Therefore the definition of \(b_2^*\) gives
\[
\int g_2(q)\,d\nu(q)=b_2^*.
\]
In particular, the definition of the Bernstein moment body implies \(b_2^*\in\mathcal B_{2,1/4}\).

It remains to prove that the membership is relative-interior membership.  We give the finite-dimensional argument so that no strict-feasibility premise is hidden.  Put
\[
C=\mathcal B_{2,1/4},\qquad b=b_2^*,\qquad A=\operatorname{aff}(C),\qquad V=A-b.
\]
The declared moment body \(C\) is compact and convex.  Suppose, for contradiction, that \(b\notin\operatorname{relint}(C)\).  By the definition of relative interior there are \(y_k\in A\setminus C\) with \(\lVert y_k-b\rVert<1/k\).  Compactness gives a closest point \(x_k\in C\) to \(y_k\).  Since \(b\in C\),
\[
0<\lVert y_k-x_k\rVert\le \lVert y_k-b\rVert<\frac1k,
\]
so \(x_k\to b\).  Define
\[
u_k=\frac{y_k-x_k}{\lVert y_k-x_k\rVert}\in V.
\]
After passage to a subsequence, compactness of the unit sphere in the finite-dimensional space \(V\) gives \(u_k\to u\in V\) with \(\lVert u\rVert=1\).

For any \(x\in C\) and \(t\in[0,1]\), convexity gives \(x_k+t(x-x_k)\in C\).  Minimality of \(x_k\) therefore implies
\[
\lVert y_k-x_k\rVert^2
\le
\lVert y_k-x_k-t(x-x_k)\rVert^2.
\]
Expanding the square, dividing by positive \(t\), and letting \(t\downarrow0\) gives
\[
u_k^{\mathsf T}(x-x_k)\le0.
\]
Taking \(k\to\infty\) along the selected subsequence yields
\[
u^{\mathsf T}(x-b)\le0\qquad\text{for every }x\in C.
\]
Apply this inequality to \(x=g_2(q)\), which belongs to \(C\) by the definition of the moment body.  Since \(b=\int g_2\,d\nu\), linearity of integration gives
\[
\int_{K_{1/4}}u^{\mathsf T}(g_2(q)-b)\,d\nu(q)
=u^{\mathsf T}\left(\int_{K_{1/4}}g_2(q)\,d\nu(q)-b\right)
=0.
\]
The integrand is continuous and nonpositive everywhere.  Full support of \(\nu\) forces it to vanish everywhere: if it were strictly negative at some point, continuity would make it strictly negative on a nonempty relatively open set, which has positive \(\nu\)-measure, contradicting the zero integral.  Consequently
\[
u^{\mathsf T}(g_2(q)-b)=0\qquad\text{for every }q\in K_{1/4}.
\]
By affine linearity the same equality holds for every \(x\in C\), hence for every \(x\in A=\operatorname{aff}(C)\).  Since \(A=b+V\) and \(u\in V\), taking \(x=b+u\in A\) gives \(u^{\mathsf T}u=0\), contrary to \(\lVert u\rVert=1\).  The contradiction proves
\[
b_2^*\in\operatorname{relint}(\mathcal B_{2,1/4}),
\]
as claimed.